\chapter{A Génese Dependente IV:\@ A Sensação Condiciona o Desejo}

\lettrine{N}{o início da prática} do \emph{paṭiccasamuppāda} está
\emph{avijjāpaccayā saṅkhārā:} a ignorância condiciona as formações kármicas.
\emph{Avijjā} é a ignorância em qualquer ser que não compreende que existe
sofrimento, o surgimento do sofrimento, a cessação do sofrimento e o Caminho que
conduz à sua cessação. Por outro lado, a palavra para conhecimento neste sentido
é \emph{vijjā}. \emph{Vijjā} é o conhecimento das Quatro Nobres Verdades: a
compreensão do sofrimento, da origem, da cessação e do Caminho.

Quando não temos insight sobre a Verdade, \emph{avijjā,} o não saber, condiciona
o \emph{saṅkhāra.} Criamos um `eu sou'. O \emph{saṅkhāra} `eu sou' é criado e
condicionado a partir desse \emph{avijjā}. Se repararem, a Primeira Nobre Verdade não
diz `Eu sofro'; a Primeira Nobre Verdade diz: `Existe sofrimento; existe
\emph{dukkha}.' Não está a dizer que alguém sofre. No entanto, pensamos que
sofremos, não é verdade? Pensamos: `Eu sofro muito nesta vida\ldots{} Ele é um
verdadeiro sofredor\ldots{} Ela sofre constantemente\ldots{} Eu sofri muito na
minha vida. Não nasci com as melhores formações kármicas disponíveis neste
planeta e tive mesmo de sofrer. Pobre de mim!' Mas o sofrimento é o que criamos
a partir da ignorância. E, por isso, o ponto importante que o Buda salientou foi
que devemos viver de acordo com o conhecimento, em vez de com a ignorância.

Esta prática budista é um caminho do conhecimento, do saber; trata-se
inteiramente de conhecer a verdade. É por isso que não sinto particularmente
pena de ninguém quando pensam que sofrem muito. Poderia dizer: `Pobrezinho.
Sinto realmente pena de ti por teres tido de sofrer.' Mas este pensamento de que
estás a sofrer não é a posição do conhecimento. Aconteceram coisas no passado,
talvez acontecimentos infelizes, e depois pensamos e nos entregamos a isso -- o
que o prolonga no presente com todo o tipo de sofrimento adicional. Mas quando
há conhecimento, insight, \emph{vijjā}, então percebemos que não há ninguém para
sofrer. Vemos as coisas como elas são.

Todo ser humano tem a capacidade de ver claramente como as coisas são e não
criar sofrimento a respeito disso.

É certo que todos nós já passámos por coisas infelizes ou fizemos coisas tolas.
Esta é apenas a experiência humana comum, não é? Quando nascemos, tudo nos pode
acontecer. Toda a gama de experiências da vida, desde as mais afortunadas às
mais infelizes, é possível para nós. Esse é o resultado do nascimento. Não há
nada de errado nisso; é apenas a forma como são as coisas. Nascer no reino
humano é arriscado -- não podemos ter a certeza do que nos espera. Pode ser uma
verdadeira confusão ou pode ser uma delícia; pode também ser, por vezes, confuso
e, por vezes, delicioso, ou um quarto confuso, três quartos medíocre e nada de
coisas deliciosas.

Nascer neste mundo humano com consciência sensorial é assim: é instável,
incerto, muda, e não podemos encontrar qualquer segurança nele. É isto que todos
temos em comum. Desde os seres humanos mais afortunados aos menos afortunados,
somos todos vulneráveis, estando numa forma e condição que podem ser
danificadas, feridas e doentes. Quando olhamos para este lado da nossa
existência humana, então não sentimos os preconceitos e as opiniões fortes sobre
raça, classe, sexo, nacionalidade e assim por diante. Somos todos irmãos e irmãs
na velhice, na doença e na morte.

Tendo nascido, existe a consciência \emph{viññāṇa}, existe o corpo,
\emph{nāma-rūpa}. Existem os órgãos dos sentidos -- \emph{saḷāyatana} -- o
olho, o ouvido, o nariz, a língua, o corpo, a mente. Existe \emph{phassa}, ou
contacto com os objetos sensoriais; e existe \emph{vedanā}, a sensação. Esta
\emph{vedanā} é o resultado do nascimento e da consciência e, neste sentido,
aplica-se à experiência sensorial das qualidades atraentes, neutras e
desagradáveis. A experiência da \emph{vedanā} através dos olhos não significa
que os seus olhos doem ou sofram; significa que, quando vê flores bonitas como
algo atraente, a \emph{vedanā} da atração é agradável. Há também sensações
desagradáveis ou neutras. Então, todo esse processo estimulará o desejo, o apego
e o devir (taṇhā-upadana-bhava). Tornamo-nos aquilo que desejamos. Agora
aplique isso a todos os sentidos e seus objetos -- ao som, ao cheiro, ao sabor,
ao toque e ao pensamento. Alguns dos nossos pensamentos são muito agradáveis;
alguns não são nem agradáveis nem desagradáveis, e alguns são desagradáveis.

Esta é a sensibilidade destes corpos; são condições totalmente sensíveis; são
conscientes e sentem. É assim que as coisas são. Alguns de vocês só querem ser
parcialmente sensíveis, não é verdade? Têm medo de serem totalmente sensíveis.
Gostariam de se tornar sensíveis apenas às coisas agradáveis, e gostariam de
rezar a Deus e dizer: `Ó Deus, por favor, dá-me tudo o que é bom, apenas
sentimentos agradáveis, e por favor torna tudo belo para mim. Nunca me deixes
sofrer, e deixa-me ter sempre sucesso, felicidade e pessoas belas à minha volta
até eu morrer\ldots{}' e essa é a mente humana chorona que quer apenas
sensibilidade parcial.

O \emph{vijjā,} ou conhecimento intuitivo, consiste em reconhecer o agradável, o
neutro e o desagradável tal como são. Já não procuramos uma sensibilidade
parcial nem as melhores experiências sensoriais, mas abrimos-nos a uma
sensibilidade total que inclui todas as possibilidades de dor, fealdade e
desagrado. \emph{Avijjā} diz: `não quero perder a minha beleza; não quero ter
nenhuma experiência desagradável; quero ser feliz.' Isso é \emph{avijjā}.
\emph{Vijjā} diz: `Existe sofrimento; existe a origem e a cessação, e existe o
caminho para sair do sofrimento.'

Contemple, então, este «eu sou» durante este retiro, este «eu sou» que chora,
lamenta-se, teme e deseja. Por que temos medo? Do que temos medo e com o que
estamos ansiosos? É a possibilidade da dor, não é? De sermos fisicamente
feridos, doentes, emocionalmente explorados ou magoados de alguma forma; de
sermos rejeitados, de não sermos amados, de sermos menosprezados, de ter cancro
ou a doença de Parkinson\ldots{} «Não quero isso; quero saúde perfeita\ldots{}
Tenho medo de poder ter alguma doença terrível\ldots{} E se tiver um daqueles
ataques cardíacos em que, durante os próximos trinta a quarenta anos, fico numa
espécie de vegetativo e os monges têm de fazer tudo, colocar-me no
penico\ldots{}? Não quero isso, não suportaria ser um incómodo ou um fardo para
ninguém.» «Não quero ser um fardo -- isso é uma obsessão inglesa, não é?»

Portanto, o «eu sou» é algo a contemplar e a observar, porque é algo que estamos
convencidos de que é a realidade para nós. Para a maioria dos seres humanos, «eu
sou» é verdade por causa da ignorância. E então é muito natural querer
felicidade e querer fugir da dor. Vês algo bonito, agarras-te a isso,
desejas-o. Algo feio -- queres livrar-te disso. Essa é a reação natural no
plano sensorial. Se for só isso, então só tens de tentar obter o melhor que
puderes e fugir de tudo o que é mau, e não há como escapar disso. É cada um por
si -- sobrevivência. Os inteligentes e os fortes sobrevivem, e os estúpidos e os
fracos ficarão no fundo, no abismo.

Mas o ser humano está dotado de uma mente reflexiva; podemos reflectir e
contemplar \emph{o vedanā.} Podemos observar e contemplar o que é a atração e o
que é a beleza. Não somos apenas animais estúpidos: podemos, de facto,
observar-nos a nós próprios a querer agarrar e possuir o belo. Podemos observar
e refletir sobre a nossa aversão a tudo o que é feio e desagradável; e também
podemos contemplar o que não é nem agradável nem desagradável.

A nossa respiração normal não é nem agradável nem desagradável; não é nem
atraente nem pouco atraente. É por isso que têm de prestar atenção a ela, porque
se a respiração fosse atraente, atrairia-vos. Eu não teria de dizer: «Observem
a vossa respiração» -- estariam a observar a vossa respiração porque era tão
atraente!

A respiração é a função fisiológica mais importante, e o corpo faz-lo quer
estejamos conscientes disso, quer sejamos loucos ou sãos, jovens ou velhos,
homens ou mulheres, ricos, pobres ou o que for. A respiração é assim. Não é
excitante nem interessante, nem repugnante ou repulsiva. Mas, à medida que nos
concentramos, que voltamos a nossa atenção para a respiração do corpo -- o que
acontece? Bem, quando me concentro na minha respiração, a mente fica tranquila;
sinto-me tranquilizado por ser capaz de me concentrar na respiração deste
corpo.

\emph{Ānāpānassati} é tédio para a maioria das pessoas no início; só inspirar --
expirar, sempre a mesma coisa. A respiração do corpo é \emph{vedanā} neutra.
Quando fazemos a meditação a percorrer as sensações do corpo, a pressão do corpo
sentado nos assentos e as roupas a tocar na pele: isso é sensação neutra.
Depois, podemos observar a \emph{vedanā} pelo ouvido, nariz, língua, olho, corpo
e mente. E começamos a ver que isso é apenas o reino sensorial, que não é uma
pessoa, que as coisas são assim apenas. Não há nada de errado, nada de mau
nisso, de todo. \emph{Vedanā} está bem. Há apenas o agradável, o doloroso e o
neutro; são apenas o que são.

No entanto, estar consciente do prazer, da dor e da \emph{vedanā} neutra,
significa termos de a suportar, de a aceitar verdadeiramente em vez de lhe
reagirmos. Refletimos sobre isso e contemplamo-la para a compreendermos
verdadeiramente. Se não contemplarmos e não tivermos insight sobre
\emph{vedanā,} continuaremos simplesmente neste processo de
\emph{paṭiccasamuppāda} -- e temos desejos, porque vedanā condiciona
\emph{taṇhā}: desejo. Mas com insight, podemos realmente quebrar o hábito.
Podemos contemplar \emph{vedanā.} Então começamos a compreender como surge o
desejo: querer o prazer, não querer a dor e simplesmente ignorar o neutro.

Uma pessoa que leva uma vida muito agitada, vai de uma coisa excitante e
estimulante, para a seguinte. Quando pensamos em estilos de vida realmente
excitantes, o que é que isso normalmente envolve? Está normalmente repleto de
tentativas frenéticas para obter experiências sensuais fantásticas; de estar
sempre em movimento -- porque a fantástica experiência sensual de ontem já
esmoreceu. Há uma necessidade de ter novas experiências sensoriais, novos
romances e aventuras, porque tudo se torna enfadonho ao se repetir. Assim,
\emph{o saṁsāra} é o ciclo, a interminável correria à procura da próxima coisa
interessante, da próxima emoção, do próximo romance, da próxima aventura -- a
próxima, a próxima, a próxima\ldots{} reparem como isso é insidioso nas nossas
vidas. Mesmo na vida monástica, mesmo num retiro de meditação, ainda podemos ser
apanhados a tentar avançar para a próxima coisa; sentados aqui a pensar no que
faremos depois do retiro, ou a tentar encontrar algo para tornar as nossas vidas
mais interessantes aqui em Amaravati.

O que é o interesse? As coisas que são interessantes são atraentes e prendem a
nossa atenção. Queremos coisas atraentes, experiências prazerosas, objetos
bonitos, música e sons bonitos. São interessantes, prendem a nossa atenção,
agradam-nos e fascinam-nos. E se uma experiência é desagradável, tememos-a.
Pode ser um inferno para a maioria das pessoas, a ideia de ter de estar num
lugar onde não há nada de belo: pessoas sombrias e enfadonhas; odores
repugnantes, grosseiros e desagradáveis; homens e mulheres sem cultura; brutos
repugnantes, imundos e nojentos; dor, doença\ldots{} É isso que tememos, aquilo
com que podemos acabar. Pode acontecer que fiquemos presos num lugar miserável.
Por isso, queremos evitar e livrar-nos de tudo isso e, depois, tentar agarrar
as experiências agradáveis tanto quanto possível.

E, no entanto, a maior parte das nossas vidas não é nem \emph{vedanā} agradável
nem dolorosa. Quando contemplam a maior parte da vossa vida, tenho a certeza de
que, para a maioria de vós, cerca de 98\% não foi nem agradável nem dolorosa,
mas apenas o que é. E, no entanto, esses 98\% da vida de uma pessoa podem passar
totalmente despercebidos porque estamos tão apegados aos extremos de esperar
pela próxima coisa, ansiando, esperando e desejando, e depois temendo e
alimentando essas possibilidades de não ter mais prazer, de não nos divertirmos.
Bem, basta pensar no nosso dia aqui em Amaravati ou em qualquer lugar do mundo.
Quanto disso é realmente prazeroso ou doloroso?

O Buda aconselhou-nos a centrar a nossa atenção nas coisas da vida que não são
nem agradáveis nem dolorosas, porque aceitar e perceber tanto o prazer como a
dor implica que tenhamos de estar atentos e alertas. Se algo não é atraente nem
repulsivo, não nos faz reagir. Não estimula de todo a nossa mente. Por isso,
temos de voltar a nossa atenção para isso, estar atentos a isso. É por isso que,
na meditação, nos sentamos, ficamos de pé, caminhamos, deitamo-nos; quatro
posturas básicas, respiração normal, coisas que são tão comuns, mas que não são
nem prazerosas nem dolorosas. A prática da atenção plena consiste em voltar a
nossa atenção para \emph{o vedanā}. Mas também não nos apegamos à neutralidade:
não estamos a tentar apegar-nos nem ao prazer nem à dor. Portanto, para estudar
\emph{vedanā,} não estamos a tentar viver uma existência neutra. Mas voltar a
atenção para ela significa que temos de nos esforçar para simplesmente sentar,
ficar de pé, caminhar, deitar-nos; estar atentos, estar aqui e agora. Temos de
prestar atenção, temos de aprender a concentrar a mente.

\emph{O vedanā} condiciona \emph{o taṇhā}. Então, o que é \emph{o taṇhā}? Esta
palavra é traduzida como desejo. É quando não estamos conscientes e alertas para
a forma como as coisas são -- então queremos ou não queremos.

Partindo do \emph{vedanā}, se for prazeroso, desejas-o; se for doloroso, não o
desejas. Depois, há um desejo sensual -- \emph{kāma-taṇhā,} o desejo de
tornar-se, a ambição. «Quero tornar-me em algo. Quero tornar-me
bem-sucedido; quero tornar-me iluminado; quero tornar-me bom. Quero
tornar-me admirado e esperado.» E \emph{vibhava-taṇhā}: desejo de me livrar de
-- esse também é forte. «Vamos livrar-nos de todas as coisas desagradáveis, dos
maus pensamentos, dos maus sentimentos, da dor, das imperfeições» = o desejo de
me livrar de.

Podemos observar estes três tipos de desejos. Podemos observá-los e refletir
sobre eles porque são objetos da mente; são objetos mentais, não são o sujeito.
O desejo não é você, por outras palavras. Mas torna-se um sujeito; torna-se
você por descuido, por \emph{avijjā.} Você agarra-se ao desejo e
\emph{torna-se} o desejado\ldots{} «Quero isto e não quero aquilo. Quero ser
bem-sucedido, não quero ser um fracasso. Tenho de me livrar destas falhas.»
Então, há o apego ao desejo, e depois tornas-te alguém que quer coisas ou não
quer coisas. E isso é interminável, não é? Quando nos \emph{tornamos} uma pessoa
que quer coisas e não quer coisas, então isso continua sem fim. Há sempre algo
que queremos e algo que não queremos. Se não vigiarmos e observarmos este
processo, então toda a nossa vida é apenas este ciclo, este ciclo interminável
de \emph{samaras} a girar e girar, apenas a desejar; a tornar-nos alguém que
quer algo, a tornar-nos alguém que não quer algo. E isso, claro, condiciona o
renascimento, \emph{jati.} Condiciona a velhice, a doença, a morte, a tristeza,
o lamento, a dor, o sofrimento e o desespero -- a depressão, a miséria;
\emph{«jara-maranam -- soka-parideva-dukkha-domanassa upayasa».}

Ser alguém que tem sempre de obter algo ou livrar-se de algo é uma forma tão
dolorosa de viver. Contemple simplesmente: qual é o verdadeiro sofrimento na sua
vida? Quando pensa que sofreu, de que é que sofreu? É de ser alguém que quer
coisas ou que não quer coisas.

Falamos da Primeira Nobre Verdade, \emph{dukkha.} Todos nós temos este
sofrimento. Quando há \emph{avijjā}, então sofremos, a nossa vida vai ser um
reino de sofrimento.

Isto está a tornar-se muito óbvio na Europa Ocidental próspera, em lugares como
a América e a Austrália, sociedades prósperas onde as pessoas conseguem muito do
que querem e onde o sofrimento não é o sofrimento da fome, da privação e da
brutalidade. Mas há tanta miséria e sofrimento nos países prósperos. De onde vem
isso? Do querer e do não querer: porque mesmo quando temos tudo o que queremos,
há mais coisas que queremos e há coisas que não queremos.

Tentar apenas satisfazer todos os nossos desejos e obter tudo o que queremos não
é a resposta, pois não? Essa não é a saída do sofrimento, porque esse processo
não termina até que o vejamos, até que usemos \emph{vijjā} em vez de
\emph{avijjā.} Contemple isso, este querer e não querer; o desejo e o apego ao
desejo.

Quando contemplas \emph{vedanā,} então vês que isso é apenas uma forma natural:
a atração e a repulsa, e nem atraente nem repulsivo. É apenas ser sensível. Por
exemplo, estas flores à minha frente são atraentes para mim neste momento. É
apenas a forma natural das coisas. Não há desejo nisso. Se eu apenas contemplar
neste momento: «Não quero essas flores», não há desejo; não quero livrar-me
delas. Não há querer ou não querer, mas elas continuam a ser agradáveis; a sua
atratividade é assim. Isso é o \emph{vedanā.} Outro exemplo é algo feio, como
estas cortinas. Acho-as feias. Sempre que entro nesta sala, a minha mente diz:
«Essas cortinas são feias.» Por isso, não se quer realmente olhar para elas.
Agora posso estar consciente do desagrado quando os meus olhos contactam essas
cortinas sem desejar livrar-me delas; é apenas a consciência da sua falta de
atratividade -- ou quando vejo a parede, que não é nem atraente nem pouco
atraente, apenas uma parede neutra.

Agora, refletindo desta forma, vê-se que é apenas a forma natural das coisas:
atração e aversão, nem atraente nem aversivo, apenas o \emph{vedanā.} Então, o
desejo é o que acrescentamos, como por aquelas flores: «Oh, quero mesmo aquelas
flores, quero ter aquelas flores no meu quarto, tenho de ter aquelas flores!»
Também as cortinas: «Quem me dera que se livrassem dessas cortinas --
incomodam-me imenso!» Ficamos obcecados em querer livrar-nos das cortinas, em
querer agarrar as flores e, claro, nem sequer reparamos na parede, a menos que
algo atraente ou pouco atraente apareça nela. E quanto ao espaço no quarto? O
espaço não é nem atraente nem pouco atraente, pois não?

Contemple, então, desta forma. O que é o desejo? Quando sente dor no corpo, se
refletir sobre a sensação física real da dor, torna-se consciente de que a essa
sensação física se junta o desejo de se livrar dela. Observe a sensação real que
tem no corpo e a aversão a ela, o desejo de se livrar da dor. Observe que a
respiração não desperta desejo. Talvez tenha o desejo de concentrar a sua mente,
de se tornar alguém que tem \emph{samādhi} ou algo do género: «Quero tornar-me
uma pessoa capaz de atingir \emph{jhāna}.»

Mas a respiração em si não é nem atraente, nem interessante, nem repulsiva. Para
a maioria das pessoas, a ideia de atingir \emph{jhāna} é atraente; tornar-se
alguém capaz de alcançar \emph{jhāna} é atraente. Assim, podemos praticar
\emph{ānāpānassati} com esse desejo; ou talvez tenhas uma mente distraída -- a
mente vagueia, não faz o que desejas. Quer que ela se concentre na respiração,
mas sempre que começa, ela vagueia. E então quer livrar-se da mente distraída,
quer tornar-se alguém que tem uma mente serena e concentrada, e não ser alguém
que tem uma mente vagueante e distraída. Portanto, há \emph{vibhava-taṇhā,} o
desejo de se livrar da mente vagueante e distraída, tornando-se alguém que tem
uma mente concentrada e pode atingir \emph{jhānas}.

Esta é uma forma de refletir sobre o desejo. Desejo pelo prazer sensorial,
desejo de tornar-se, desejo de se livrar. Se realmente contemplarmos e
conhecermos \emph{vedanā} exatamente através de \emph{vijjā}, através da atenção
plena e da sabedoria, então não criamos desejo. Ainda há o prazer, a dor, o nem
prazer nem dor, mas as coisas são como são. Esta é a \emph{verdadeira natureza},
a forma como as coisas são; é o Dhamma, a Verdade. Portanto, não há sofrimento
quando as coisas são como são. O sofrimento é o resultado de
desejo-apego-devir (\emph{taṇhā-upadana-bhava}). A partir daí, a sequência
de \emph{paṭiccasamuppāda} leva ao nascimento, envelhecimento, morte, tristeza,
lamentação, dor, pesar e desespero
(\emph{jati-jara-mananam-soka-parideva-dukkha-domanassa-upayasa}). Toda a
sequência de sofrimento decorre de \emph{taṇhā-upadana-bhava.}

Portanto, contemple este tema do \emph{paṭiccasamuppāda} durante este retiro. O
desejo de se livrar do desejo continua a ser uma armadilha da mente, não é
verdade? A contemplação não é livrar-se, mas compreender. Este é o caminho do
conhecimento, do \emph{vijjā} em vez do \emph{avijjā.}

